\section{Which kernel occurs in the comparison of periods?}

% BEGIN S89K1
\paragraph{S89K1. The two restrictions at the last pair of lifts.}
In the last pair of lifts there are two orthogonal decompositions
$V_0=A\perp L$ and $X=T\perp L'$, with $L,L'$ lines. The group
$U(S)$ occurs in both lifts. Fix the standard additive character in the source's sign convention and
the Weil splittings of the two pairs. On the two line factors the
splitting characters are $\chi$ and $\chi^{-1}$, and on the orthogonal
sums the splitting characters are the corresponding products. The
line map $\tau:L'\to L$ satisfies
\[
       \tau^*H_L\tau=-\xi_0\bar\xi_0 H_{L'} .
\]
Here $H$ denotes the Gram matrix and $\xi_0$ is the chosen imaginary
generator. In orthogonal coordinates write $\Xi(u,r)$ for the first
theta test on $(V_0,S)$ and $\Phi(v)$ for the second on $(S,T)$.
The common two-variable datum is explicitly
\[
 C\bigl(u;(v,w)\bigr)=\Phi(v)\,\Xi(u,\kappa w),\qquad
 \kappa w=-\xi_0^{-1}H_S^{-1}(\tau w)^*H_L .
\]
Here $w$ is a matrix for the $(S,L')$ variable, and $u,r,v$ are the
variables for $(A,S),(L,S),(S,T)$, respectively. All Gram matrices
in this coordinate formula are extended to the adeles.

For an element $g$ of the first group of a pair and $h$ of the
second, write $[\omega(g,h)f]_0$ for evaluation of the transformed
test at the origin. For $a\in U(A)(\mathbb A)$,
$h\in U(S)(\mathbb A)$ and $b\in U(T)(\mathbb A)$ one has
\[
\begin{split}
 [\omega_{S,T}(h,b)\Phi]_0\,
 [\omega_{V_0,S}(a\oplus 1,h)\Xi]_0
 ={}&
 \bigl[\omega_{A,S}(a,h)\bigl(
 u\longmapsto
 [\omega_{S,X}(h,b\oplus 1) C(u;\cdot)]_0
 \bigr)\bigr]_0 .
\end{split}
\]
To see this, fix $u$ and put $f_u(r)=\Xi(u,r)$. In the $X$ decomposition
the test is the tensor $\Phi\otimes(f_u\circ\kappa)$. The tensor
rule splits its transform into the $T$ and $L'$ factors. The change
of variables $\kappa$ on the line intertwines the action of $h$
on $L'$ with its action on $L$. Indeed, multiplying $\tau$ by
$\overline{-\xi_0^{-1}}$ gives an isometry from the negative of
$L'$ to $L$ by the displayed Gram equation. Passing to the adjoint
model interchanges source and target and conjugates the Weil action.
Negating the form and inverting its line character conjugates it
again. The two conjugations cancel, giving
$\omega_{S,L'}(h,1)(f_u\circ\kappa)
 =(\omega_{L,S}(1,h)f_u)\circ\kappa$ in the specified splittings.
Thus the inner value on the right is
$[\omega_{S,T}(h,b)\Phi]_0[\omega_{L,S}(1,h)f_u]_0$.
Take out the first scalar and apply the tensor rule in
$A\perp L$ to the second factor as a function of $u$. This gives
the left side, including its scalar. This is the pointwise kernel
identity used before integration. Kudla's \emph{Seesaw dual reductive pairs} (1984), \S1,
(1.14), explains why two pairings of this restriction coincide when
the integrals can be interchanged.
% END S89K1

% BEGIN S89K2
\paragraph{S89K2. One kernel serves every tuple at a fixed pair of steps.}
Hold fixed the spaces, orthogonal decompositions, two common theta
tests, and the link on the line. At each interior contraction also fix
its additive domain and measure. Use the continuous theta tests with
the prescribed Heisenberg transformation law. The assignments used
by individual tuples for the passage to the smaller Weil model give
the same reduced test. Indeed this reduction takes Fourier
coefficients in the two isotropic blocks of a Witt decomposition.
In each integration variable a rational translation multiplies the
theta test by its Heisenberg character. The factor in the Fourier
coefficient supplies its inverse, since the translated block is
isotropic and orthogonal to the middle block. The integrand thus
descends to the additive adelic quotient. Its absolute value is periodic and bounded on the compact additive
quotient, so these Fourier integrals converge absolutely. Integration uses
Haar measure of total mass one, so either choice of fundamental
domain and covolume-one measure gives the same value. The remaining
unipotent averaging measure is fixed by its prescribed right-invariant
normalization. This proves equality of the reduced tests, including
the normalizing scalar, in both directions of the passage.

For an interior pair let $f_1$ and $f_2$ be the two reduced tests.
In the orthogonal coordinates on the two ends the common datum is
their contraction
\[
                C(u,v)=\int_{\mathcal F_L}
                     f_2(v,w)\,f_1(u,\kappa_j w)\,dw .
\]
Here $\mathcal F_L$ and $dw$ are the common additive matrix
fundamental domain and measure for the line being contracted, and
$\kappa_j$ is the link on that line. The preceding independence
applies separately to $f_1,f_2$. For the terminal pair, $\Xi$ in
S89K1 is the first reduced test and $\Phi$ is the restriction of
the second test to the middle space. This proves tuple-independence
in that case as well. Applying the same level operator to either
common datum preserves equality. In particular the kernel may be
constructed using any one tuple when the family is nonempty.
To identify integration against that kernel with the prescribed
period in each tuple one still needs the finite-place unfolding, including
the support at the level and twist places, the values on the
surviving cells and the volume constants. Those local computations are hypotheses of the family period
transfer used below.
% END S89K2

% BEGIN S89K3
\paragraph{S89K3. The relation transported by the chain is a one-dimensional factorization.}
Fix a word and an even chain of length $m=2d\ge2$. For one tuple
$t$, a permitted bottom column $\phi$, and a terminal row $F$, let
$A_q$ be the terminal scalar period after lifting the column up the
chain, and $D_q(y,v)$ the bottom Bessel period after lifting the row
down, where $q=(t,\phi,F)$. Here $y$ is a bottom frame parameter
and $v$ a smooth row test. Include in $A_q$ the factor
\[
 c_P(t)\prod_{i=0}^{d-1}
 c_V(t,2i)c_K(t,2i)c_K(t,2i+1),
\]
and in $D_q$ the factor $\prod_{j=0}^{m-1}c_I(t,j)$. These are,
respectively, the terminal evaluation and two-step column constants,
and the normalizations of the successive row lifts.

Suppose the terminal and interior period transfer laws hold, with
the permitted column spaces stable under the successive column
operators. There is a single function $H(y,v)$, valid for all such
tuples, columns and rows, such that
\[
                         D_q(y,v)=A_q H(y,v).
\]
Here is the argument. Take any finite relation
$\sum_q\alpha_q A_q=0$. The terminal transfer gives the zero
relation of tested periods two levels below. Each interior transfer
carries this zero relation down two more levels, for every frame
and smooth test there. At a pair starting at $j$, the three factors
$c_V(t,j)c_K(t,j)c_K(t,j+1)$ on the column side are exchanged
for the two factors $c_I(t,j)c_I(t,j+1)$ on the row side. The
terminal transfer also uses $c_P(t)$. This proves
$\sum_q\alpha_q D_q=0$ with precisely the products displayed above. For $d=1$ the terminal transfer already gives this conclusion.

If some $A_{q_0}\ne0$, take the relation with coefficients
$A_{q_0}$ at $q$ and $-A_q$ at $q_0$. It gives
$H=D_{q_0}/A_{q_0}$, independent of the choice of $q_0$. If all
$A_q$ vanish, the one-term relations give $D_q=0$ and $H=0$ works.
Conversely the factorization immediately implies transport of every
finite relation. This is the exact content of the source's
\texttt{MainX}. On the family route its hypotheses follow from the
two-step unfolding and common frame expansions at the chosen
places, together with the terminal evaluation, cusp decay and the
smooth-vector spanning argument. The factorization allows each word
its own $H$. To identify an arithmetic period with a Hilbert-space
pairing one must also match the terminal scalar and the lower
period to those two readings on the selected data and supply the
nonzero comparison value there.
% END S89K3
